% LuaLaTeX Document - Generated by Media Scribe Workflow
\documentclass[a4paper,10pt,twocolumn]{ltjsarticle}

% LuaLaTeX用フォント設定パッケージ
\usepackage{luatexja-fontspec}
\usepackage{amsmath,amssymb}
\usepackage{unicode-math}

% ====================
% 欧文フォント設定 (Libertinus)
% ====================
\setmainfont{Libertinus Serif}[
    BoldFont = {Libertinus Serif Bold},
    ItalicFont = {Libertinus Serif Italic},
    BoldItalicFont = {Libertinus Serif Bold Italic}
]
\setsansfont{Libertinus Sans}[
    BoldFont = {Libertinus Sans Bold},
    ItalicFont = {Libertinus Sans Italic}
]
\setmonofont{DejaVu Sans Mono}[Scale=0.9]

% ====================
% 日本語フォント設定 (原ノ味フォント)
% ====================
\setmainjfont{HaranoAjiMincho-Regular}[
    BoldFont = {HaranoAjiGothic-Medium},
    ItalicFont = {HaranoAjiMincho-Regular},
    BoldItalicFont = {HaranoAjiGothic-Bold}
]
\setsansjfont{HaranoAjiGothic-Regular}[
    BoldFont = {HaranoAjiGothic-Bold}
]
\setmonojfont{HaranoAjiGothic-Regular}

% ====================
% 数式フォント設定 (Libertinus Math)
% ====================
\setmathfont{Libertinus Math}

% ファイル生成日時（JST）- 生成時に置換
\newcommand{\generatedDate}{2026-01-12}
\newcommand{\generatedTime}{23:38}

% その他のパッケージ
\usepackage[margin=20mm]{geometry}
\usepackage{hyperref}
\usepackage{booktabs}
\usepackage{array}
\usepackage{tabularx}
\usepackage{enumitem}
\usepackage{graphicx}
\usepackage{ascmac}

% ハイパーリンクの色設定
\hypersetup{
    colorlinks=true,
    linkcolor=blue,
    urlcolor=blue,
    citecolor=blue
}

% ヘッダー・フッター設定
\usepackage{fancyhdr}
\usepackage{lastpage}
\pagestyle{fancy}
\fancyhf{}
\fancyhead[R]{\small \generatedDate\ \generatedTime\ JST (\thepage/\pageref{LastPage})}
\renewcommand{\headrulewidth}{0.4pt}

% 1ページ目のスタイル（ページ番号なし）
\fancypagestyle{firstpage}{
    \fancyhf{}
    \fancyhead[R]{\small \generatedDate\ \generatedTime\ JST}
    \renewcommand{\headrulewidth}{0.4pt}
}

% Y列タイプ定義（tabularx用）
\newcolumntype{Y}{>{\raggedright\arraybackslash}X}

\begin{document}
\thispagestyle{firstpage}

\title{ホルンレッスン: アンブシュア --- 講師: 濵地 宗（群馬交響楽団 首席ホルン奏者） --- 2026-01-15}
\author{ましノート\\創価大学 新世紀管弦楽団}
\date{}

\maketitle

\tableofcontents
\newpage

\section{レッスン概要}

\begin{itemize}
  \item 日付: 2026-01-15
  \item 講師: 濵地 宗（群馬交響楽団 首席ホルン奏者）
  \item 楽器: ホルン
  \item テーマ: アンブシュア
\end{itemize}

\section{レッスン内容}

\texttt{[00:00:00.000]} はい、今日はアンブシュアについてレッスンをしていきたいと思います。

\texttt{[00:00:06.000]} アンブシュアというのは、マウスピースに当てる唇の形のことですね。

\texttt{[00:00:12.500]} ホルンの場合は、上唇と下唇のバランスが非常に重要になってきます。

\texttt{[00:00:18.500]} まず、基本的な形としては、口角を少し引いて、
中央に向かって息を集めるイメージです。

\texttt{[00:00:25.500]} よくある間違いとして、唇を強く押し付けすぎることがあります。

\texttt{[00:00:32.500]} これは、音が出にくくなったり、
音色が硬くなる原因になります。

\texttt{[00:00:38.500]} では、実際に吹いてみましょう。
リラックスして、自然な形で唇をマウスピースに当てます。

\texttt{[00:00:45.500]} そう、そんな感じです。
もう少し上唇を使うイメージで。

\texttt{[00:00:52.500]} 良いですね。
今の感覚を覚えておいてください。

\texttt{[00:00:58.500]} 次に、高音域に行くときのアンブシュアの変化について説明します。


\section{テクニック解説}

% techniques セクション - 内容を追加してください


\section{練習課題}

% exercises セクション - 内容を追加してください


\section{用語集}

% 専門用語の解説をここに追加
\textit{（用語集を追加してください）}

\section{練習のポイント}

% self_practice セクション - 内容を追加してください


\section{Summary}

% AI分析によるサマリーをここに追加
\textit{（サマリーを追加してください）}

\section*{謝辞}

本レポートは Claude Code を使用して生成されました。

\end{document}